\newcounter{english}
\documentclass{article}

% packages
\usepackage{amsmath, amsthm, thmtools, amsfonts, amssymb, luacode, catchfile, hyperref, ifthen, graphicx, bbm}
\ifcsname c@kobocompile\endcsname
	\usepackage[a5paper, total={1072pt, 1448pt}, margin=10pt, includeheadfoot]{geometry} % set page margins
\else
	\usepackage[a4paper, margin=50pt, includeheadfoot]{geometry}
\fi
\usepackage[shortlabels]{enumitem}
\usepackage[skip=3pt, indent=0pt]{parskip}

% language
\usepackage[bidi=basic, layout=tabular, provide=*]{babel}
\ifcsname c@english\endcsname
	\babelprovide[main, import]{english}
\else
	\babelprovide[main, import, alph=letters.gershayim]{hebrew}
	\setlist[enumerate,2]{label={\alph*.}}
	\babelprovide{rl}
\fi
\babelfont{rm}[Renderer=Harfbuzz]{Libertinus Serif}
\babelfont{sf}{Libertinus Sans}
\babelfont{tt}{Libertinus Mono}

% style
\ifcsname c@stoppagebreaks\endcsname
\else
	\AddToHook{cmd/section/before}{\clearpage}	% Add line break before section
\fi
\linespread{1.3}
\setcounter{secnumdepth}{0}		% Remove default number tags from sections, this won't do well with theorems
\AtBeginDocument{\setlength{\belowdisplayskip}{3pt}}
\AtBeginDocument{\setlength{\abovedisplayskip}{3pt}}
\graphicspath{ {../images/} }

% operators
\DeclareMathOperator\cis{cis}
\DeclareMathOperator\Sp{Span}
\DeclareMathOperator\tr{tr}
\DeclareMathOperator\im{Im}
\DeclareMathOperator\re{Re}
\DeclareMathOperator\diag{diag}
\DeclareMathOperator*\lowlim{\underline{lim}}
\DeclareMathOperator*\uplim{\overline{lim}}
\DeclareMathOperator\rng{rng}
\DeclareMathOperator\Sym{Sym}
\DeclareMathOperator\Arg{Arg}
\DeclareMathOperator\Log{Log}
\DeclareMathOperator\dom{dom}
\DeclareMathOperator\supp{Supp}
\DeclareMathOperator\var{Var}
\DeclareMathOperator\cov{Cov}
\DeclareMathOperator\End{End}
\DeclareMathOperator\id{id}
\DeclareMathOperator\cl{cl}
\DeclareMathOperator\gal{Gal}
\DeclareMathOperator\dist{dist}
\DeclareMathOperator\form{form}
\DeclareMathOperator\term{term}
\DeclareMathOperator\sent{sent}
\DeclareMathOperator\Th{Th}
\DeclareMathOperator\Hom{Hom}
\DeclareMathOperator\Aut{Aut}
\DeclareMathOperator\Char{char}
\DeclareMathOperator\vol{vol}
\DeclareMathOperator\GL{GL}
\DeclareMathOperator\SL{SL}
\DeclareMathOperator\lcm{lcm}

% commands
\newcommand{\cupdot}{\mathbin{\mathaccent\cdot\cup}}
\newcommand{\indicator}{\mathbbm{1}}
\newcommand{\AAA}[0]{\mathbb{A}}
\newcommand{\Aa}[0]{\mathcal{A}}
\newcommand{\BB}[0]{\mathbb{B}}
\newcommand{\Bb}[0]{\mathcal{B}}
\newcommand{\CC}[0]{\mathbb{C}}
\newcommand{\Cc}[0]{\mathcal{C}}
\newcommand{\DD}[0]{\mathbb{D}}
\newcommand{\Dd}[0]{\mathcal{D}}
\newcommand{\EE}[0]{\mathbb{E}}
\newcommand{\FF}[0]{\mathbb{F}}
\newcommand{\Ff}[0]{\mathcal{F}}
\newcommand{\Ii}[0]{\mathcal{I}}
\newcommand{\Gg}[0]{\mathcal{G}}
\newcommand{\HH}[0]{\mathbb{H}}
\newcommand{\Hh}[0]{\mathcal{H}}
\newcommand{\Kk}[0]{\mathcal{K}}
\newcommand{\Ll}[0]{\mathcal{L}}
\newcommand{\Mm}[0]{\mathcal{M}}
\newcommand{\NN}[0]{\mathbb{N}}
\newcommand{\Nn}[0]{\mathcal{N}}
\newcommand{\PP}[0]{\mathbb{P}}
\newcommand{\Pp}[0]{\mathcal{P}}
\newcommand{\QQ}[0]{\mathbb{Q}}
\newcommand{\RR}[0]{\mathbb{R}}
\newcommand{\Rr}[0]{\mathcal{R}}
\newcommand{\Ss}[0]{\mathcal{S}}
\newcommand{\TT}[0]{\mathbb{T}}
\newcommand{\Tt}[0]{\mathcal{T}}
\newcommand{\UU}[0]{\mathbb{U}}
\newcommand{\Uu}[0]{\mathcal{U}}
\newcommand{\Vv}[0]{\mathcal{V}}
\newcommand{\Ww}[0]{\mathcal{W}}
\newcommand{\Xx}[0]{\mathcal{X}}
\newcommand{\Yy}[0]{\mathcal{Y}}
\newcommand{\ZZ}[0]{\mathbb{Z}}
\newcommand{\Zz}[0]{\mathcal{Z}}
\newcommand{\acts}[0]{\circlearrowright}
\newcommand{\explain}[2] {
	\begin{flalign*}
		 && \text{#2} && \text{#1}
	\end{flalign*}
}
\newcommand{\maketitleprint}[1][green]{
	\begin{center}
		\includegraphics[width=6cm]{dragon_#1}
	\end{center}
}

% theorem commands
\newtheoremstyle{c_remark}
	{}	% Space above
	{}	% Space below
	{}% Body font
	{}	% Indent amount
	{\bfseries}	% Theorem head font
	{}	% Punctuation after theorem head
	{.5em}	% Space after theorem head
	{\thmname{#1}\thmnumber{ #2}\thmnote{ \normalfont{\text{(#3)}}}}	% head content
\newtheoremstyle{c_definition}
	{3pt}	% Space above
	{3pt}	% Space below
	{}% Body font
	{}	% Indent amount
	{\bfseries}	% Theorem head font
	{}	% Punctuation after theorem head
	{.5em}	% Space after theorem head
	{\thmname{#1}\thmnumber{ #2}\thmnote{ \normalfont{\text{(#3)}}}}	% head content
\newtheoremstyle{c_plain}
	{3pt}	% Space above
	{3pt}	% Space below
	{\itshape}% Body font
	{}	% Indent amount
	{\bfseries}	% Theorem head font
	{}	% Punctuation after theorem head
	{.5em}	% Space after theorem head
	{\thmname{#1}\thmnumber{ #2}\thmnote{ \text{(#3)}}}	% head content

\ifcsname c@english\endcsname
	\theoremstyle{plain}
	\newtheorem{theorem}{Theorem}[section]
	\newtheorem*{theorem*}{Theorem}
	\newtheorem{lemma}[theorem]{Lemma}
	\newtheorem{proposition}[theorem]{Proposition}
	\newtheorem*{proposition*}{Proposition}

	\theoremstyle{definition}
	\newtheorem{definition}[theorem]{Definition}
	\newtheorem*{definition*}{Definition}
	\newtheorem{example}{Example}[section]
	\newtheorem{exercise}{Exercise}[section]

	\theoremstyle{remark}
	\newtheorem*{remark}{Remark}
	\newtheorem*{solution}{Solution}
	\newtheorem{corollary}[theorem]{Corollary}
	\newtheorem{notation}[theorem]{Notation}
\else
	\theoremstyle{c_plain}
	\newtheorem{theorem}{משפט}[section]
	\newtheorem*{theorem*}{משפט}
	\newtheorem{lemma}[theorem]{למה}
	\newtheorem{proposition}[theorem]{טענה}
	\newtheorem*{proposition*}{טענה}

	\theoremstyle{c_definition}
	\newtheorem{definition}[theorem]{הגדרה}
	\newtheorem*{definition*}{הגדרה}
	\newtheorem{example}{דוגמה}[section]
	\newtheorem{exercise}{תרגיל}[section]

	\theoremstyle{c_remark}
	\newtheorem*{remark}{הערה}
	\newtheorem*{solution}{פתרון}
	\newtheorem{corollary}[theorem]{מסקנה}
	\newtheorem{notation}[theorem]{סימון}
\fi

% Questions related commands
\newcounter{question}
\setcounter{question}{1}
\newcounter{sub_question}
\setcounter{sub_question}{1}

\ifcsname c@english\endcsname
	\newcommand{\question}[1][0]{
		\ifthenelse{#1 = 0}{}{\setcounter{question}{#1}}
		\section{Question \arabic{question}}
		\addtocounter{question}{1}
		\setcounter{sub_question}{1}
	}

	\newcommand{\subquestion}[1][0]{
		\ifthenelse{#1 = 0}{}{\setcounter{sub_question}{#1}}
		\subsection{Part \alph{sub_question}}
		\addtocounter{sub_question}{1}
	}
\else
	\newcommand{\question}[1][0]{
		\ifthenelse{#1 = 0}{}{\setcounter{question}{#1}}
		\section{שאלה \arabic{question}}
		\addtocounter{question}{1}
		\setcounter{sub_question}{1}
	}

	\newcommand{\subquestion}[1][0]{
		\ifthenelse{#1 = 0}{}{\setcounter{sub_question}{#1}}
		\subsection{סעיף \localecounter{letters.gershayim}{sub_question}}
		\addtocounter{sub_question}{1}
	}
\fi

% import lua and start of document
\directlua{common = require ('../common')}

\GetEnv{AUTHOR}

% headers
\author{\AUTHOR}
\date\today

\title{Lie Groups --- Journals}
\setcounter{secnumdepth}{2}
% chktex-file 9
% chktex-file 17

\usepackage{fancyhdr}
\pagestyle{fancy}
\renewcommand{\headrulewidth}{0pt}

\DeclareMathOperator{\SU}{SU}

\begin{document}
\maketitle
\maketitleprint[teal]{}

\section{Journal 1}
We will discuss several definitions of a smooth manifold on the special case that the manifold $X$ is a euclidean manifold (embedded manifold), namely $X \subseteq \RR^n$ for some $n \in \NN$.
Note that $\RR^n$ is a T4 topological space of 2nd countability, therefore if $A$ is an atlas on $X$ then $(X, A)$, or for the sake of convenience $A$ alone, is a smooth manifold.
When discussing embedded manifolds, there are several definitions, they are generally give substance to the concept of a manifold.

The first definition is Guillemin \& Pollacks (GP) definition, as shown in their book ``Differential Topology''.
\begin{definition}[GP smooth manifold]
	Let $M \subseteq \RR^n$ be a set and $k \in \NN$.
	We say that $M$ is a $k$-dimensional smooth manifold if that for any point $p \in M$ there is an open neighborhood $p \in U \subseteq \RR^n$ and an open set $W \subseteq \RR^k$ with function $F : U \to W$ diffeomorphism,
	meaning that its smooth and invertible. We define the function $f = F|_{M \cap U}$ as a smooth function from open subset of the manifold $M$.
\end{definition}
Another version which we consider is defined by its parameterizations,
\begin{definition}[Locally smooth parameterization manifold]
	a set $M \subseteq \RR^n$ together with a $k \in \NN$ is called a smooth manifold if for any $p \in M$ there is an open set $U \subseteq M$ (open in the sub-topology of $\RR^n$ induced by $M$),
	and an open set $W \subseteq \RR^k$ such that there exists a function $\varphi : W \to U$ smooth (in the standard sense), open, invertible and with $\rank D_x = k$ for any $x \in W$, namely a smooth local parameterization.
\end{definition}
While both definitions may seem different, it can be shown using the inverse function theorem that they are equivalent.

Whitney embedding theorem states that any smooth manifold $(X, A)$ can be embedded in $\RR^{2 k}$ space where $k = \dim (A, X)$,
giving us motivation to compare our definition of a manifold in euclidean space to the definitions stated above.
\begin{proposition}
	If $M \subseteq \RR^n$ is a GP $k$-dimensional smooth manifold and for any $p \in M$ let $f_p : U_p \to W_p$ be a diffeomorphism, where $U_p \subseteq M, W_p \subseteq \RR^k$ open.
	Then $A = \{ f_p \mid p \in M \}$ is an atlas, and in particular $(M, A)$ is a smooth manifold.
\end{proposition}
\begin{proof}
	Clear by our definition of $f_p$ as a restriction of a smooth function $F_p$ and compositions.
\end{proof}
This result together with Whitney embedding theorem lead us to the conclusion that a manifold is a smooth manifold iff its a GP manifold (up to embedding).

\section{Journal 2}
We will discuss the claim that $\exp : L(G) \to G$ is a local homeomorphism, and in particular the fact that $\ker \exp$ is a discrete set.
In general, the following theorem is known about embedded manifolds,
\begin{theorem}[Preimage theorem]
	If $y$ is a regular value of $f : X \to Y$, then $f^{-1}(Y)$ is a submanifold of $X$ with $\dim f^{-1}(y) = \dim X - \dim Y$.
\end{theorem}
where,
\begin{definition}[Regular value]
	If $f : X \to Y$ is a smooth map between manifolds, then $y \in Y$ is called a regular value if $d f_x : T_x(X) \to T_y(Y)$ is surjective at every point $x \in f^{-1}(y)$.
\end{definition}
If a function is locally homeomorphism, then in particular it is locally surjective, and defines a preimage manifold with $\dim f^{-1}|_U(y) = \dim X - \dim T_e X = 0$, namely it is a discrete set of points.
While this is a very complex way to perceive this result, it binds the behavior of the exponential map to intersection-theory, and open the question about these two topics.
Let us proceed to prove the claim in discussion directly.
\begin{proposition}
	If $G$ has a Lie group structure, and $\exp : L(M) \to M$ is the exponential function, then $\exp^{-1}(g)$ is discrete for any $g \in G$.
\end{proposition}
\begin{proof}
	If $g \notin G_e$ then $\exp^{-1}(g) = \emptyset$, then let us assume that $\exp^{-1}(g)$ is non-empty, in particular $g \in G_e$.
	Assume for the sake of contradiction that a point $x \in \exp^{-1}(g)$ is non-discrete, namely there exist a sequence $\{ x_n \} \subseteq \exp^{-1}(g)$ with $x_n \to x$.
	$\exp$ is a diffeomorphism, then let $\varphi : \overset{\subseteq G}{U} \to \overset{\subseteq L(G)}{W}$ be a chart, then $h = \varphi^{-1} \circ \exp$ is a diffeomorphism,
	but $\exp(x_n) = g$ therefore $h(x_n)$ is constant for any $n$, a contradiction.
\end{proof}

\section{Journal 3}
We will bring a full and formal proof of the closed subgroup theorem, assuming that the three lemmas Adams formulate hold.
\begin{theorem}[Closed subgroup]
	If $G$ is a lie group and $H \le G$ is a closed subgroup, then $H$ assumes a smooth manifold structure.
\end{theorem}
\begin{proof}
	Denote $V$ the vector space for which the atlas maps $G$ locally.
	By the third lemma there is a vector space $W \le V$ such that $\exp W \subseteq H$ is a neighborhood of $0$, in particular there is an open neighborhood of $e$ in $H$, denoted by $U$.
	$\exp$ is a local diffeomorphism, by restricting $\exp$ to $W$ or a smaller set $\Omega$, we can assume that $\exp \Omega = U$ and $\exp |_{\Omega} : \Omega \to U$ is invertible.
	Note that to achieve $\Omega \subseteq W, U \subseteq \exp W$ we can use the inverse map theorem.
	Define $\varphi : U \to \Omega$ as the inverse of the restriction of $\exp$, we get a function of a neighborhood of $e$ in $H$.
	For any $p \in H$, define $\varphi_p : U_p \to \Omega$. Recall that the group multiplication is smooth, therefore the preimage of an open neighborhood of $p$, small enough, for the map $x \mapsto x \cdot p$,
	yields an open set such that $U_p = p U$ (where we assume that $U$ is already small enough without loss of generality).
	Then $\varphi_p(x) = \varphi(p^{-1} x)$ is exactly composition of smooth maps therefore smooth, and is a chart around $p$.
	Construct $A = \{ \varphi_p : U_p \to \Omega \mid p \in H \}$, a collection of smooth maps such that $\bigcup_{p \in H} \dom \varphi_p = H$, therefore $(H, A)$ is indeed a smooth map.
\end{proof}

\section{Journal 4}
We will discuss examples of Lie groups.
In class we saw that if $G$ is an Abelian Lie group, then it is isomorphic to $T^n \times \RR^k$ where $T^n$ is an $n$-torus.
We also know that $\GL_n(\FF)$ is a Lie group for $\FF \in \{\RR, \CC, \HH\}$ and $n \in \NN$.
Lastly, it can be shown that $O(n), U(n)$ are Lie groups as closed subgroups of $\GL_n$.
\begin{proposition}
	$\SU(n) = \{ M \in U(n) \mid \det M = 1 \}$, the group of unitary complex matrices of order $n \times n$ with determinant $1$, forms a Lie group structure.
\end{proposition}
\begin{proof}
	$U(n)$ has a Lie group structure, therefore it suffices to show that $\SU(n) \le U(n)$ is a closed subgroup.
	It is clear that it is a subgroup, as $I_n \in \SU(n)$, and the fact that determinant is multiplicative.
	Let $f : U(n) \to \RR$ be the function defined by $f(M) = \det M$, then $f$ is continuous as a polynomial, therefore $f^{-1}(\{ 1 \})$ is closed as the preimage of the closed subset $\{ 1 \}$.
	But this preimage is exactly $\SU(n)$, therefore by the closed subgroup theorem, it assumes the structure of a submanifold, hence has a Lie group structure.
\end{proof}
\begin{proposition}
	As smooth manifolds, $\SU(3) \cong S^3$, namely there is a smooth bijective map between the manifolds defined for these structures.
\end{proposition}
\begin{proof}
	We identify $S^3$ as,
	\[
		S^3
		= \{ (z, w) \in \CC \times \CC \mid {|z|}^2 + {|w|}^2 = 1 \}
	\]
	and define,
	\[
		\varphi : S^3 \to \SU(3),
		\qquad
		\varphi(z, w)
		= \begin{pmatrix}
			z & -\bar{w} \\
			w & \bar{z}
		\end{pmatrix}
	\]
	As a function over $\RR^4$ it is linear, and therefore injective iff $\ker \varphi = \{ 0 \}$, by direct check it indeed is injective.
	It is also surjective, therefore a bijection.
	The proof that it also preserves smoothness of charts is clear and is not present to save space.
\end{proof}
Geometrically, the result we achieved is that $S^3$ is an embedded Lie group, thus providing us with another example of a concrete, somewhat visualization-able.

As a bonus we want to describe Hopf-fibration.
Define Hopf map, $h : S^3(1) \to S^2(\frac{1}{2})$ by
\[
	h(z, w)
	= (z \overline{w}, \frac{1}{2} ({|z|}^2 - {|w|}^2))
\]
We show that its image is subset of $S^2(\frac{1}{2})$ and that any $q \in S^2(\frac{1}{2})$ is a regular value of $h$.
\begin{proof}
	Notice that $\frac{1}{2} ({|z|}^2 - {|w|}^2) \in \RR$ and $z \overline{w} \in \CC$, then $h$ is embedded in $\RR^3$.
	Also,
	\[
		{|h(z, w)|}^2
		= z \overline{w} w \overline{z} + \frac{1}{4} {|z|}^4 - \frac{1}{2} {|zw|}^2 + \frac{1}{4} {|w|}^4
		= \frac{1}{4} {|z|}^4 + \frac{1}{2} {|zw|}^2 + \frac{1}{4} {|w|}^4
		= {\left(\frac{1}{2} ({|z|}^2 + {|w|}^2)\right)}^2
		= {\left(\frac{1}{2} \cdot 1\right)}^2
	\]
	It follows that $\im h \subseteq S^2(\frac{1}{2})$.

	As for regularity (By Guillemin \& Pollacks definition) notice that,
	\[
		D h |_{(z, w)}
		= D h |_{(z_r, z_i, w_r, w_i)}
		= \begin{pmatrix}
			w_r & - w_i & z_r & -z_i \\
			w_i & w_r & z_i & z_r \\
			\frac{1}{2} z_r & \frac{1}{2} z_i & \frac{1}{2} w_r & \frac{1}{2} w_i
		\end{pmatrix} 
		= \begin{pmatrix}
			\overline{w} & \overline{z} \\
			i \overline{w} & i \overline{z} \\
			\frac{1}{2} z & \frac{1}{2} w
		\end{pmatrix} 
	\]
	then the point $(z, w)$ is regular if the map is onto, $\RR^2 \simeq \CC \simeq T_{h(z, w)}(S^2(\frac{1}{2}))$.
	By linear independence of the matrix columns each $(z, w) \ne 0$ is a regular value, but $0 \notin S^3(1)$, therefore any $q \in S^2(\frac{1}{2})$ is regular value.
\end{proof}
By looking into intersection theory of embedded smooth manifolds we can deduce that $S^1 \to S^2 \xrightarrow{h} S^3$ is a fibration, in particular it assumes Lie group structure.

\section{Journal 5}
We will complete the second part of Schur's lemma,
\begin{lemma}[Schur]
	Let $G$ be a topological group.
	\begin{enumerate}
		\item If $V$ and $W$ are irreducible representations of $G$ over $\FF$,
			then any $\FF G$-map $f : V \to W$ is either an isomorphism or $0$.
		\item If $V$ is an irreducible $G$-representation over $\CC$, then $\End_{\CC G}(V) \cong \CC$.
	\end{enumerate}
\end{lemma}
\begin{proof}
	The first part was proved during the lecture, and us achieved through the fact that $\ker f \subseteq V$ and $\im f \subseteq W$ are both sub-representations.
	Together with the fact that there are no proper non-trivial sub-representations to both $V$ and $W$, either $\ker f = \{ 0 \}$ or $\im f = \{ 0 \}$.

	For the second part, let $T \in \End_{\CC G}(V)$ be some endomorphism.
	As a linear operator over a finite dimensional vector space over $\CC$, there is some eigenvalue $\lambda \in \CC$, meaning that there is $v \in V^\times$ with $T v = \lambda v$,
	or equivalently $(T - \lambda I) v = 0$, where $I$ is identity map in $\End(V)$. Then $T - \lambda I$ has a non-trivial kernel, therefore by the first part of the lemma $T - \lambda I \equiv 0$.
	Therefore $T \equiv \lambda I$. Define,
	\[
		\varphi : \CC \to \End_{\CC G}(V),
		\qquad
		\varphi(\lambda)
		= \lambda I
	\]
	we get a homomorphism such that it is surjective (by the previous statement) and injective, therefore $\End_{\CC G}(V) \cong \CC$.
\end{proof}

\end{document}
